% =========================================================================
% Capítulo 3: Cinética de sistemas de partículas
% Tema 01: Sistemas de partículas
% Problema 05: Mancuerna deslizante en pista circular
% Origen: Ayudantía 28, 2026
% =========================================================================
\subsubsection{Mancuerna deslizante en pista circular}

Una mancuerna está constituida por dos partículas $A$ y $B$, cada una de masa $m$, conectadas mediante una barra rígida sin masa de longitud $2l$. Los extremos de la mancuerna se deslizan sin fricción por el interior de una pista circular fija de radio $R$ en el plano vertical. Sea $\theta$ el ángulo que forma la vertical descendente con la línea que une el centro de la pista $O$ con el centro de la barra $C$. Considerando que la gravedad actúa hacia abajo, realice lo siguiente:

\begin{enumerate}
    \item Defina los marcos de referencia y sistemas de coordenadas adecuados.
    \item Analice la cinemática del sistema, estableciendo la relación geométrica entre la distancia $a$ (distancia $OC$) y las constantes $R$ y $l$.
    \item Analice la cinética del sistema para determinar la ecuación diferencial del movimiento en términos del ángulo $\theta$.
\end{enumerate}

\begin{figure}[!ht]
    \centering
    \begin{tikzpicture}[>=Latex, scale=1.5]
        % --- Parámetros ---
        \def\R{2.5}       % Radio de la pista
        \def\halfL{1.5}   % Mitad de la longitud de la barra (l)
        % Por Pitágoras, la distancia a = sqrt(R^2 - l^2). 
        % Con R=2.5 y l=1.5 -> a = sqrt(6.25 - 2.25) = 2.0
        \def\distA{2.0}   
        \def\ang{30}      % Ángulo theta

        % Coordenada Origen
        \coordinate (O) at (0,0);

        % --- 1. Pista Circular ---
        % Pared exterior (hachurada)
        \draw[thick] (O) circle (\R+0.2);
        \begin{scope}
            \clip (O) circle (\R+0.4);
            \fill[pattern=north east lines, even odd rule] (O) circle (\R+0.2) circle (\R+0.4);
        \end{scope}
        \draw[thick] (O) circle (\R+0.2); % Redibujar borde

        % --- 2. Elementos Rotados (Barra y Radio) ---
        \begin{scope}[rotate=-\ang] % Rotación horaria theta
            
            % Línea OC (distancia a)
            \draw[dashed] (O) -- (0, -\distA) node[midway, left] {$a$};
            \coordinate (C) at (0, -\distA);

            % Barra AB (Perpendicular a OC)
            \draw[line width=4pt] (-\halfL, -\distA) -- (\halfL, -\distA);
            \draw[line width=3pt, gray!40] (-\halfL, -\distA) -- (\halfL, -\distA);
            \draw[thick] (-\halfL, -\distA) -- (\halfL, -\distA);

            % Partículas
            \coordinate (A) at (-\halfL, -\distA);
            \coordinate (B) at (\halfL, -\distA);

            \shadedraw[ball color=gray!30] (A) circle (0.2) node[right, xshift=5pt, yshift=5pt] {$A$};
            \shadedraw[ball color=gray!30] (B) circle (0.2) node[right, xshift=5pt] {$B$};
            \node[font=\footnotesize] at (A) {$m$}; % Opcional
            \node[font=\footnotesize] at (B) {$m$};

            % Radio R (hacia B)
            \draw[->, thin] (O) -- (\R+0.2,0) node[midway, above right] {$R$};
            
            % Cota 2l
            \draw[|<->|] ($(A)+(0,-0.3)$) -- ($(B)+(0,-0.3)$) node[midway, below] {$2l$};
            \filldraw (C) circle (0.6pt) node[above right, xshift=2pt, font=\footnotesize] {$C$};
        \end{scope}

        % --- 3. Referencias y Ángulos ---
        % Eje Vertical
        \draw[dashed, gray] (O) -- (0, -2.5);
        
        % Ángulo theta
        \draw[->] (0, -1.0) arc (-90:{-90-\ang}:1.0);
        \node at (-0.3, -1.3) {$\theta$};

        % Centro O
        \filldraw (O) circle (1.5pt) node[above left] {$O$};

        % Gravedad
        \draw[->, thick] (3.0, 0) -- (3.0, -1.0) node[midway, right] {$\mathbf{g}$};

    \end{tikzpicture}
    \caption{Sistema de mancuerna deslizante dentro de una pista circular vertical.}
    \label{fig:mancuerna_circular}
\end{figure}

% --- SOLUCIÓN ---------------------------------------------------------------
\begin{solucion}{Mancuerna deslizante en pista circular}
Dado que no hay fricción y el sistema está sujeto únicamente a la fuerza de gravedad (que es conservativa) y a las fuerzas normales (que no realizan trabajo al ser perpendiculares a la trayectoria), la energía mecánica total del sistema se conserva ($\frac{dE}{dt} = 0$). Este es el enfoque más directo para obtener la ecuación de movimiento.

\subsubsection*{Geometría del sistema}
El triángulo formado por el origen $O$ y las partículas $A$ y $B$ es isósceles, ya que los lados $OA$ y $OB$ equivalen al radio $R$. El centro de la varilla es el punto $C$, y la distancia desde $A$ a $C$ es $l$.

\begin{center}
\begin{tikzpicture}[>=Latex, scale=1.4]
    % --- PARÁMETROS ---
    \def\R{2.5}    % radio de la pista [u]
    \def\l{1.5}    % mitad de la varilla [u]
    \def\a{2.0}    % distancia OC [u]

    % --- GEOMETRÍA ---
    \coordinate (O) at (0,0);
    \coordinate (C) at (0,-\a);
    \coordinate (A) at (-\l,-\a);
    \coordinate (B) at (\l,-\a);
    \draw[thick] (A) -- (B);
    \draw[dashed] (O) -- (C) node[midway, left] {$a$};
    \draw (O) -- (A) node[midway, above left] {$R$};
    \draw (O) -- (B) node[midway, above right] {$R$};
    \draw (-0.2,-\a) -- (-0.2,-\a+0.2) -- (0,-\a+0.2);

    % --- PUNTOS ---
    \foreach \P/\pos in {A/left, B/right}
        \shadedraw[ball color=gray!30] (\P) circle (0.13) node[\pos=4pt, yshift=-3pt] {$\P$};
    \filldraw (O) circle (1.2pt) node[above] {$O$};
    \filldraw (C) circle (1.2pt) node[below] {$C$};
    \draw[|<->|] ($(A)+(0,-0.45)$) -- ($(C)+(0,-0.45)$) node[midway, below] {$l$};
\end{tikzpicture}
\end{center}

La distancia perpendicular desde el centro $O$ a la varilla $AB$ (segmento $OC$) se denota por $a$. Aplicando el teorema de Pitágoras en el triángulo rectángulo $OCA$:
$$a = \sqrt{R^2 - l^2}$$
Como la varilla es rígida y sus extremos siempre tocan el círculo, la distancia $OC = a$ permanece constante durante todo el movimiento. El conjunto entero (las dos masas y la varilla) rota de manera rígida alrededor del punto $O$.

\subsubsection*{Energía del sistema}
\textbf{Energía cinética ($T$)}\\
El sistema gira como un cuerpo rígido alrededor de $O$ con una velocidad angular $\dot{\theta}$. Como las partículas $A$ y $B$ se encuentran siempre a una distancia $R$ del origen, la magnitud de la velocidad tangencial de ambas es $v_A = v_B = R\dot{\theta}$. La energía cinética total es la suma de las energías cinéticas de ambas partículas:
$$T = \frac{1}{2}m v_A^2 + \frac{1}{2}m v_B^2 = \frac{1}{2}m(R\dot{\theta})^2 + \frac{1}{2}m(R\dot{\theta})^2 = m R^2 \dot{\theta}^2$$

\textbf{Energía potencial gravitacional ($V$)}\\
Fijamos el nivel de referencia gravitacional ($V = 0$) en el eje horizontal que pasa por $O$, y consideramos el eje vertical hacia abajo como positivo. La posición del centro de masa del sistema compuesto por las dos partículas idénticas se ubica exactamente en el punto $C$. La coordenada vertical del centro de masa es $y_C = a\cos\theta$. La energía potencial de ambas masas en conjunto es equivalente a concentrar toda la masa ($2m$) en su centro de masa $C$:
$$V = -(2m)g y_C = -2mg a\cos\theta = -2mg\sqrt{R^2 - l^2}\cos\theta$$

\subsubsection*{Conservación de la energía y ecuación de movimiento}
La energía mecánica total del sistema es $E = T + V$:
$$E = m R^2 \dot{\theta}^2 - 2mg\sqrt{R^2 - l^2}\cos\theta = \text{constante}$$
Para encontrar la ecuación diferencial de movimiento, derivamos $E$ con respecto al tiempo $t$ y la igualamos a cero:
$$\frac{dE}{dt} = 2m R^2 \dot{\theta} \ddot{\theta} - 2mg\sqrt{R^2 - l^2}(-\sin\theta)\dot{\theta} = 0$$
$$2m\dot{\theta} \left( R^2\ddot{\theta} + g\sqrt{R^2 - l^2}\sin\theta \right) = 0$$
Para un movimiento general ($\dot{\theta} \neq 0$), el término entre paréntesis debe ser nulo. Al dividir todo entre $R^2$, obtenemos la ecuación diferencial que describe el movimiento de la mancuerna:
\respuesta{$\displaystyle \ddot{\theta} + \frac{g\sqrt{R^2 - l^2}}{R^2}\sin\theta = 0$}

\patotruco{Fíjate que las normales de la pista no aparecen en ningún momento: son perpendiculares a la velocidad y no trabajan. Cuando el problema pide solo la ecuación de movimiento y no las reacciones, la energía te ahorra proyectar fuerzas.}

\patotip[squareglasses]{Comprobación rápida: si $l\to0$ la mancuerna colapsa en una partícula doble de masa $2m$ que cae por una pista de radio $R$, y la ecuación se reduce a la del péndulo simple $\ddot\theta+\frac{g}{R}\sin\theta=0$.}
\end{solucion}
